\documentclass[../../main.tex]{subfiles}
\begin{document}
\section{对 Lie 理论的应用}\label{section:ch21-applications-lie-theory}

本节收集商流形定理与齐性空间的一些推论。它们给出 Lie 群上光滑映射与同态的基本事实，最终结束第 8 章提出的关于连通 Lie 群由 Lie 代数刻画的项目。

首先是关于抽象群的两个引理。

\begin{theorem}{[抽象群商定理]}\label{theorem:ch21-抽象群商}
设 $G$ 是群，$H\subseteq G$ 是正规子群，$G/H$ 为左陪集集合。则 $G/H$ 关于 $(g_1H)(g_2H)=(g_1g_2)H$ 的乘法成为群，且商映射 $\pi\colon G\to G/H$ 是同态，核为 $H$。
\end{theorem}

\begin{proof}
$H$ 正规意味着 $gH=Hg$（$g\in G$），这保证该乘法良定义：若 $g_1H=g_1^\ast H$、$g_2H=g_2^\ast H$，则 $g_1g_2H=g_1g_2HH=g_1H g_2H=g_1^\ast H g_2^\ast H=g_1^\ast g_2^\ast H$。群公理由 $G$ 继承。$\pi$ 的核为 $H$ 显然。
\end{proof}

\begin{theorem}{[抽象群第一同构定理]}\label{theorem:ch21-抽象群第一同构}
设 $F\colon G\to H$ 是群同态。则 $F$ 的核是 $G$ 的正规子群，像为 $H$ 的子群，且由 $\widetilde F(g\ker F)=F(g)$ 给出的映射 $\widetilde F\colon G/\ker F\to H$ 是单同态，其像为 $F(G)$。
\end{theorem}

\begin{proof}
核正规、像子群由同态定义易得。$\widetilde F$ 良定义单射：$F(g_1)=F(g_2)$ 当且仅当 $g_1^{-1}g_2\in\ker F$。
\end{proof}

\begin{exercise}\label{exercise:ch21-抽象群定理证明}
证明上述两定理。
\end{exercise}

\begin{theorem}{[Lie 群商定理]}\label{theorem:ch21-Lie群商}
设 $G$ 是 Lie 群，$H\subseteq G$ 是正规闭 Lie 子群。则 $G/H$ 关于陪集乘法是 Lie 群，商映射 $\pi\colon G\to G/H$ 是满射光滑同态，其李代数为 $\mathfrak g/\mathfrak h$。
\end{theorem}

\begin{proof}
由齐性空间构造定理，$G/H$ 是光滑流形，$\pi$ 是光滑淹没。乘法 $\mu\colon G/H\times G/H\to G/H$ 良定义光滑（用定理 4.30，纤维上常值），逆亦然。故 $G/H$ 是 Lie 群，$\pi$ 是光滑同态。
\end{proof}

\begin{theorem}{[Lie 群第一同构定理]}\label{theorem:ch21-Lie群第一同构}
设 $F\colon G\to H$ 是 Lie 群同态。则 $\ker F$ 是 $G$ 的闭正规 Lie 子群，$F(G)$ 是 $H$ 的 Lie 子群，且诱导映射 $\widetilde F\colon G/\ker F\to H$ 是单射光滑浸入；若 $F$ 具有常秩，则它到像上是光滑淹没。
\end{theorem}

\begin{proof}
核闭正规，故由闭子群定理是闭正规 Lie 子群，$G/\ker F$ 是 Lie 群。$\widetilde F$ 由秩定理光滑浸入；常秩时为淹没。
\end{proof}

\begin{proposition}\label{proposition:ch21-离散子群闭}
Lie 群的离散子群是闭的。
\end{proposition}

\begin{proof}
设 $\Gamma\subseteq G$ 离散，$(g_i)\subset\Gamma$ 收敛于 $g\in G$。则 $g_i^{-1}g_{i+1}\to e$。因 $\Gamma$ 离散，$e$ 在 $G$ 中有邻域 $U$ 使 $U\cap\Gamma=\{e\}$。$i$ 大时 $g_i^{-1}g_{i+1}\in U\cap\Gamma$，故 $g_i^{-1}g_{i+1}=e$，即 $g_i=g_{i+1}$。最终 $g_i$ 常数，故 $g\in\Gamma$。
\end{proof}

\begin{theorem}{[离散子群的商是覆叠]}\label{theorem:ch21-离散子群商覆叠}
设 $G$ 是连通 Lie 群，$\Gamma\subseteq G$ 是离散子群。则 $\Gamma$ 在 $G$ 上由左乘光滑自由真作用，商映射 $\pi\colon G\to G/\Gamma$ 是正规覆叠映射。
\end{theorem}

\begin{proof}
$\Gamma$ 闭故 Lie 子群。左乘作用自由。真性由命题\cref{proposition:ch21-离散作用真判别} 验证。由定理\cref{theorem:ch21-离散群商覆叠}，$\pi$ 是光滑正规覆叠。
\end{proof}

\begin{example}{[离散商]}\label{example:ch21-离散商}
$\mathbb{R}^n/\mathbb{Z}^n\cong\mathbb{T}^n$；$\mathbb{S}^n/\{\pm1\}\cong\mathbb{RP}^n$；$\operatorname{SU}(2)/\{\pm I\}\cong\operatorname{SO}(3)$。
\end{example}

\begin{theorem}{[离散核同态]}\label{theorem:ch21-离散核同态}
设 $F\colon G\to H$ 是连通 Lie 群间的满射光滑同态，$\ker F$ 离散。则 $F$ 是正规覆叠映射，且 $\ker F$ 同构于 $H$ 的基本群。
\end{theorem}

\begin{proof}
核离散闭，故 $F$ 的局部逆由反函数定理存在；$F$ 是覆叠。其甲板变换群即核，故 $\ker F\cong\pi_1(H)$。
\end{proof}

\begin{theorem}{[给定 Lie 代数的连通群]}\label{theorem:ch21-给定代数连通群}
每个有限维实 Lie 代数 $\mathfrak g$ 都是某连通 Lie 群 $G$ 的 Lie 代数。
\end{theorem}

\begin{proof}
由 Ado 定理 $\mathfrak g$ 同构于某 $\operatorname{GL}(n,\mathbb{R})$ 的 Lie 子代数；其连通积分李子群即所求。
\end{proof}

\begin{proposition}\label{proposition:ch21-连通性判据}
若 Lie 群 $G$ 由含单位元的连通开集生成，则 $G$ 连通。
\end{proposition}

\begin{proposition}\label{proposition:ch21-典型群连通性}
$\operatorname{SO}(n)$、$U(n)$、$\operatorname{SU}(n)$ 对所有 $n$ 连通；$O(n)$ 恰有两个连通分支 $O(n)^+$ 与 $O(n)^-$，由行列式正负区分。
\end{proposition}

\begin{proposition}\label{proposition:ch21-GL分支}
$\operatorname{GL}(n,\mathbb{R})$ 恰有两个连通分支，由行列式符号区分；$\operatorname{GL}(n,\mathbb{C})$ 连通。
\end{proposition}

\begin{figure}[ht]
\centering
\includegraphics[width=0.8\textwidth]{fig-21-2.png}
\caption{$GL^+(n,\mathbb{R})$ 连通性的证明}\label{figure:ch21-glc-connected}
\end{figure}
\end{document}
